\documentclass[conference]{IEEEtran}

\usepackage[utf8]{inputenc}
\usepackage[T1]{fontenc}
\usepackage{graphicx}
\usepackage{amsmath}
\usepackage{amssymb}
\usepackage{cite}
\usepackage{hyperref}
\usepackage{xcolor}
\usepackage{float}
\usepackage{caption}
\usepackage{siunitx}
\usepackage{booktabs}
\usepackage{tabularx}


\begin{document}

\title{Learning the Theory of Monte Carlo Methods}

\author{
    \IEEEauthorblockN{Estiven Castrillon Alzate}
    \IEEEauthorblockA{
        Instituto de Física, Universidad de Antioquia \\
        \href{mailto:estiven.castrillon2@udea.edu.co}{estiven.castrillon2@udea.edu.co}
    }
}

\maketitle

\begin{abstract}
This document presents a study of the Monte Carlo acceptance-rejection method based on the reference \textit{Probability for Physicists} by Simon Sirca. The method is applied to the differential decay rate of the muon, including radiative corrections, with the goal of understanding sampling techniques used in High Energy Physics simulations.
\end{abstract}

\section{Introduction}

The Monte-Carlo method is a tool used widely in computer simulations in experimental physics. Here, I will study this mathematical and computational technique to enhance my understanding of simulations in High Energy Physics (HEP). The main reference for this learning process is the book \textit{Probability for Physicists}~\cite{Sirca_2016}.

The Monte-Carlo method considered in this work is the acceptance-rejection method, a technique used to generate random samples from a probability distribution that may be difficult to sample from directly. The method works by generating random samples from a proposal distribution and accepting or rejecting those samples according to a criterion that guarantees that the accepted samples follow the desired target distribution~\cite{Nakade_2025}.

\subsection{Acceptance-Rejection Method}

This will be the method used in the Monte Carlo analysis developed in this repository, since the probability density function (PDF) of the decay spectrum is already known, it can be directly employed as the target distribution in the acceptance rejection sampling method, in our case, the target distribution corresponds to the differential decay rate $\frac{d\Gamma}{dx}$, where \(x\) is the normalized electron energy defined as


\[
x = \frac{2E_e}{m_\mu}, \qquad 0 < x \leq 1,
\]


with \(E_e\) the electron energy and \(m_\mu\) the muon mass.

Neglecting terms of order \(m_e^2/m_\mu^2\) and higher-order radiative corrections \(\mathcal{O}(\alpha^2)\), the differential decay rate for free muon decay is given by


\[
\frac{d\Gamma^{\mathrm{free}}}{dx} = \frac{G_F^2 m_\mu^5}{192\pi^3} \,x^2 \left( 6 - 4x + \frac{\alpha}{\pi} f(x) \right),
\]


where \(G_F\) is the Fermi constant and $\alpha = \frac{1}{137.035999173(35)}$ is the fine-structure constant.

The function \(f(x)\) represents the radiative corrections to the electron energy spectrum and is explicitly written as:

\begin{equation}
\begin{aligned}
f(x) &= \Bigg[ \frac{5}{3x^2} + \frac{16x}{3} + \frac{4}{x} + (12 - 8x)\ln\!\left(\frac{1}{x}-1\right) - 8 \Bigg]\ln\!\left(\frac{m_\mu}{m_e}\right) \\
&\quad + (6 - 4x)\Bigg[ 2\,\mathrm{Li}_2(x) - 2\ln^2(x) + \ln(x) \\
&\qquad + \ln(1-x)\left(3\ln(x) - \frac{1}{x} - 1\right) - \frac{\pi^2}{3} - 2 \Bigg] \\
&\quad + \frac{(1-x)}{3x^2}\Big[ 34x^2 + (5 - 34x^2 + 17x)\ln(x) - 22x \Big] \\
&\quad + 6(1-x)\ln(x).
\end{aligned}
\end{equation}

This expression is implemented directly in the Monte Carlo framework and used as the target probability density in the acceptance rejection algorithm for generating electron energy spectra from muon decay~\cite{Czarnecki_2014}.

\section*{Comportamiento límite de las funciones en el Monte Carlo}

Sea la tasa diferencial de decaimiento


\[
\frac{d\Gamma}{dx} \equiv \_dGamma\_dx\_differential\_decay\_rate(x),
\]


la cual ya incluye el factor global $x^2$ en su definición interna.

La parte radiativa se modela mediante


\[
f(x) \equiv \_f\_radiative\_corrections(x),
\]


que contiene términos singulares como $\frac{1}{x}$, $\frac{1}{x^2}$ y $\ln(x)$.

Para analizar su contribución efectiva, se considera el producto


\[
x^2 \cdot f(x).
\]



\subsection*{Límites en los bordes}

\begin{itemize}
\item Cuando $x \to 0^+$:


\[
\lim_{x \to 0^+} x^2 f(x) \quad \text{es finito}.
\]


\item Cuando $x \to 1^-$:


\[
\lim_{x \to 1^-} x^2 f(x) \quad \text{es también finito}.
\]


\end{itemize}

\subsection*{Conclusión}

El producto $x^2 f(x)$ asegura que la parte radiativa aporta una contribución bien definida en ambos extremos del dominio $x \in (0,1)$.

\bibliographystyle{IEEEtran}
\bibliography{references}

\end{document}